% ===================================================================
%  Vorlage: Quintenzirkel (Dur außen, Moll innen, Vorzeichenzahl)
%  Reines TikZ – kein pgfplots, kein musixtex nötig.
%    pdflatex vorlage_quintenzirkel.tex
% ===================================================================
\input{header.tex}
\input{newcommands.tex}
\pagestyle{empty}

\begin{document}
\begin{center}
{\Large\bfseries\sffamily Quintenzirkel}\\[2pt]
{\small Dur-Tonarten außen, parallele Moll-Tonarten innen, Anzahl der Vorzeichen am Rand}
\end{center}
\vspace{4mm}

\begin{center}
\begin{tikzpicture}[font=\sffamily]
  \def\rdur{3.25}   % Radius Dur-Ring
  \def\rmoll{2.05}  % Radius Moll-Ring
  \def\rvz{4.35}    % Radius Vorzeichen-Beschriftung

  % pos (0 = oben, im Uhrzeigersinn) / Dur / Moll / Vorzeichen
  \foreach \pos/\dur/\moll/\vz in {%
    0/C/a/{0},
    1/G/e/{1\(\sharp\)},
    2/D/h/{2\(\sharp\)},
    3/A/fis/{3\(\sharp\)},
    4/E/cis/{4\(\sharp\)},
    5/H/gis/{5\(\sharp\)},
    6/{Fis/Ges}/{dis/es}/{6\(\sharp\)\,/\,6\(\flat\)},
    7/Des/b/{5\(\flat\)},
    8/As/f/{4\(\flat\)},
    9/Es/c/{3\(\flat\)},
    10/B/g/{2\(\flat\)},
    11/F/d/{1\(\flat\)}}
  {
    \pgfmathsetmacro{\ang}{90 - 30*\pos}
    \node[circle, draw=akzentblau!70!black, fill=akzentblau!15,
          minimum size=12mm, inner sep=0pt, line width=0.5pt]
      (D\pos) at (\ang:\rdur) {\textbf{\dur}};
    \node[circle, draw=akzentorange!70!black, fill=akzentorange!15,
          minimum size=9mm, inner sep=0pt, line width=0.5pt]
      (M\pos) at (\ang:\rmoll) {\small\moll};
    \node[text=gray!55!black] at (\ang:\rvz) {\footnotesize\vz};
  }
  \node[align=center, text=gray!50!black] at (0,0)
    {\scriptsize\textbf{Dur}\\[-1pt]\tiny außen\\[3pt]\scriptsize\textbf{Moll}\\[-1pt]\tiny innen};
\end{tikzpicture}
\end{center}

\end{document}
